\chapter{Concluzii}
Proiectul de diplomă prezentat în această lucrare a avut ca scop dezvoltarea unei platforme web pentru studiul limbii engleze care să reunească într-un singur loc materialele de studiu, evaluarea și sesiunile cu profesori, pornind de la observația că soluțiile existente acoperă de obicei doar una dintre aceste componente, obligând utilizatorul să jongleze cu mai multe servicii în același timp.

Rezultatul final este o aplicație funcțională care acoperă ciclul întreg de învățare, de la înregistrarea contului până la susținerea sesiunilor video, toate modulele principale fiind implementate și funcționale. Fluxurile de aprobare pentru profesori și cursuri, sistemul de credite, integrarea cu Stripe și Daily.co, urmărirea progresului, toate acestea funcționează conform așteptărilor stabilite la începutul proiectului. Cel mai reușit aspect al implementării este tocmai coerența sistemului, fiecare componentă legându-se firesc de celelalte în loc să existe ca module independente fără nicio legătură între ele.

Au existat și dificultăți pe parcurs, precum gestionarea rezervărilor simultane și prevenirea deducerilor de credite de mai multe ori, care au necesitat mai multă atenție, dar și procesarea webhook-urilor Stripe, care a impus un sistem de idempotență pentru a evita procesarea dublă a evenimentelor. Fluxul de activare pentru utilizatorii minori a durat, de asemenea, mai mult timp decât era planificat inițial, gestionarea stărilor intermediare ale contului dovedindu-se mai complexă decât părea.

Platforma are limitări clare care nu pot fi ignorate. Catalogul de cursuri nu are funcționalitate de filtrare, ceea ce poate deveni o problemă pe măsură ce numărul de cursuri crește. Disponibilitatea profesorilor nu permite excepții punctuale, un profesor neputând marca o zi liberă sau disponibilitate suplimentară în afara intervalelor recurente configurate.

Direcțiile de dezvoltare ulterioară sunt destul de clare în viziunea mea. Un sistem de recenzii pentru cursuri ar ajuta studenții să aleagă mai bine, iar un mecanism de recomandări bazat pe nivelul și istoricul de studiu al fiecărui student ar face experiența mult mai personalizată. Pe partea de conținut, adăugarea unor exerciții audio pentru pronunție sau exerciții de conversație ghidată ar extinde semnificativ tipurile de activități disponibile. Toate acestea rămân pentru o versiune viitoare a platformei.
